% Демонстрация 2: «Конструктор кривых 2-го порядка» (Лекция 9)
\section*{Демонстрация 2. «Конструктор кривых 2-го порядка»\\ \normalsize\textmd{Файл \texttt{demos/02-conics.html} \quad(к Лекции~9)}}

Демонстрация представляет собой интерактивный «конструктор», в котором все
конические сечения порождаются \emph{единым} полярным уравнением
\[
  r = \frac{p}{1 - e\cos\varphi},
\]
а их вид целиком определяется двумя ползунками: эксцентриситетом $e$ и
фокальным параметром $p$. Предусмотрены: переключатель формы записи
(\textbf{полярная} $\leftrightarrow$ \textbf{каноническая}), подсветка
\textbf{фокуса} (полюса) и \textbf{директрисы}, отображение координатной
сетки и осей. По мере движения ползунка $e$ программа автоматически
определяет тип кривой (окружность $\to$ эллипс $\to$ парабола $\to$
гипербола), пересчитывает параметры $a$, $b$, $c$ и показывает
соответствующее каноническое уравнение.

\subsection*{Какие факты лекции иллюстрирует}

Демонстрация наглядно связывает воедино материал раздела «Кривые второго
порядка» (определения эллипса, гиперболы и параболы и
уравнения~(9.16), (9.18), (9.19)).

\begin{itemize}
  \item \textbf{Единая природа конических сечений.} Непрерывное изменение
    эксцентриситета $e$ от $0$ до значений, больших единицы, показывает,
    что окружность, эллипс, парабола и гипербола — это не четыре
    разрозненных объекта, а \emph{одно семейство} кривых, различающихся
    лишь одним числом. При $e = 0$ имеем окружность $x^2 + y^2 = r^2$ (это
    случай $\varepsilon = 0$ из лекции, когда эллипс «превращается в
    окружность»); при $0 < e < 1$ — эллипс~(9.16); при $e = 1$ — параболу
    $y^2 = 2px$~(9.19); при $e > 1$ — гиперболу~(9.18). Ползунок физически
    «проходит» через все эти случаи без разрывов.

  \item \textbf{Геометрический смысл эксцентриситета.} В лекции
    $\varepsilon = c/a$ вводится формально, как «мера сплюснутости».
    Демонстрация делает эту меру осязаемой: при $e \to 0$ эллипс округляется
    (фокусы сходятся к центру, $c \to 0$), при $e \to 1^{-}$ эллипс
    неограниченно вытягивается, его дальняя вершина уходит на
    бесконечность — и в пределе кривая \emph{размыкается}, становясь
    параболой. Панель параметров синхронно показывает, как растут $c$ и $a$
    и как выполняется тождество $c^2 = a^2 - b^2$ (эллипс) или
    $c^2 = a^2 + b^2$ (гипербола).

  \item \textbf{Роль фокуса и директрисы.} Полюс полярной системы совпадает
    с фокусом $F$ кривой — он подсвечен красным маркером. Директриса
    (жёлтая штриховая прямая) отвечает тем самым прямым
    $x = \pm a/\varepsilon$ из лекции для эллипса и гиперболы и прямой
    $x = -p/2$ для параболы. Тем самым параметр $e$ получает второй,
    \emph{фокально-директориальный} смысл: $e$ есть отношение расстояния от
    точки кривой до фокуса к расстоянию до директрисы. Для параболы это
    отношение равно $1$ — ровно то определение параболы
    ($KM = FM$), из которого выведено уравнение~(9.19).

  \item \textbf{Связь канонической и полярной форм.} Переключатель режима
    показывает на одной и той же кривой две записи: полярную
    $r = p/(1 - e\cos\varphi)$, единую для всех типов, и каноническую,
    свою для каждого типа. Это иллюстрирует переход от полярных координат к
    декартовым ($x = r\cos\varphi$, $y = r\sin\varphi$) из раздела
    «Полярные координаты» и показывает, что фокальный параметр $p$ связан с
    полуосями соотношением $p = b^2/a$.
\end{itemize}

\begin{remarkIT}{Где это работает на практике}{demo02_apps}
  Единое полярное уравнение конических сечений — не учебная абстракция, а
  рабочий инструмент:
  \begin{itemize}
    \item \textbf{Орбитальная механика и Кеплер.} Траектория тела в поле
      тяготения — коническое сечение с фокусом в притягивающем центре, и
      записывается она именно как $r = p/(1 - e\cos\varphi)$. Эксцентриситет
      $e$ различает связанные орбиты (эллипс, $e < 1$, спутники и планеты),
      параболический уход ($e = 1$) и гиперболический пролёт ($e > 1$,
      межзвёздные тела). Расчёт эфемерид и работа спутниковых
      навигационных систем (GPS/ГЛОНАСС) опираются на эту формулу.
    \item \textbf{Антенны и оптика.} Параболоид фокусирует параллельный
      пучок в фокус (спутниковые «тарелки», телескопы-рефлекторы), а
      эллиптические и гиперболические зеркала переносят излучение из одного
      фокуса в другой — основа двухзеркальных схем Кассегрена и медицинских
      литотриптеров.
    \item \textbf{Траектории и физические движки.} Баллистические
      траектории снарядов в играх и симуляторах — параболы и эллипсы;
      управление эксцентриситетом задаёт «настильность» или навесность
      полёта.
  \end{itemize}
\end{remarkIT}

\subsection*{Примерный сценарий использования на занятии}

Демонстрацию удобно включить \emph{после} вывода канонических уравнений
эллипса, гиперболы и параболы (уравнения~(9.16), (9.18), (9.19)) — как
обобщающий мостик, связывающий три отдельных вывода в одну картину.

\begin{enumerate}
  \item \textbf{Старт от знакомого.} Откройте демонстрацию в режиме
    \emph{каноническая} с $e = 0$. На экране — окружность $x^2 + y^2 = r^2$.
    Зафиксируйте внимание: «это случай $\varepsilon = 0$, о котором мы
    говорили; фокус совпал с центром». Включите подсветку фокуса и сетки.

  \item \textbf{Рождение эллипса.} Медленно увеличивайте $e$ до $0{,}5$.
    Фокус отъезжает от центра, окружность сплющивается. Спросите:
    \emph{«Что в панели параметров отвечает за смещение фокуса?»} (ответ:
    растёт $c = ea$). Обратите внимание на выполнение тождества
    $c^2 = a^2 - b^2$ в строке параметров.

  \item \textbf{Подключаем директрису.} Включите флажок «Директриса».
    Подчеркните: положение жёлтой прямой — это $x = -a/\varepsilon$ из
    лекции. Поставьте вопрос: \emph{«Куда уйдёт директриса, если $e$
    устремить к нулю?»} (ответ: на бесконечность — у окружности директрисы
    нет).

  \item \textbf{Приближение к границе.} Доведите $e$ до $0{,}9$, затем до
    $0{,}97$. Эллипс вытягивается, его дальняя вершина уползает за край
    экрана. Спросите класс: \emph{«При каком $e$ кривая разомкнётся и
    перестанет быть замкнутой?»} Дайте студентам высказать гипотезы.

  \item \textbf{Момент удивления: парабола как граница.} Поставьте ровно
    $e = 1$. Кривая размыкается — эллипс «лопнул» и стал параболой
    $y^2 = 2px$. Это ключевой момент: парабола — это не «ещё одна кривая», а
    \emph{пороговое} состояние между эллипсом и гиперболой. Здесь же
    напомните фокально-директориальное определение параболы:
    для параболы расстояния до фокуса и до директрисы равны, то есть в
    точности $e = 1$.

  \item \textbf{За порогом — гипербола.} Сдвиньте $e$ до $1{,}5$. Появляются
    две ветви, на экране проступают асимптоты, тождество в панели меняется
    на $c^2 = a^2 + b^2$. Подчеркните смену знака в каноническом уравнении:
    $\dfrac{x^2}{a^2} - \dfrac{y^2}{b^2} = 1$ против
    $\dfrac{x^2}{a^2} + \dfrac{y^2}{b^2} = 1$ у эллипса — один знак
    отличает замкнутую кривую от разомкнутой.

  \item \textbf{Связь двух языков.} Переключите тумблер
    \emph{полярная}\,/\,\emph{каноническая} при фиксированном $e$. Покажите,
    что одна и та же кривая описывается единым полярным уравнением
    $r = p/(1 - e\cos\varphi)$ и «своим» каноническим. Задайте вопрос:
    \emph{«Почему полярная запись одна на всех, а каноническая — разная?»}
    (полюс помещён в фокус, общий для всего семейства).

  \item \textbf{Роль параметра $p$.} Меняя только ползунок $p$ при
    фиксированном $e$, покажите, что $p$ задаёт «размер» кривой, не меняя её
    типа, и связан с полуосями как $p = b^2/a$. Завершите выводом: тип
    кривой определяет \emph{только} эксцентриситет $e$, а $p$ —
    масштабирующий множитель.
\end{enumerate}

\begin{remarkIT}{Методическая подсказка}{demo02_tip}
  Не объявляйте тип кривой заранее — пусть его «угадывает» сам ползунок и
  подпись на панели. Наибольший дидактический эффект даёт именно
  \emph{переход через $e = 1$}: предложите студентам предсказать, что
  произойдёт, прежде чем выставить это значение. Возврат ползунка назад
  (гипербола $\to$ парабола $\to$ эллипс $\to$ окружность) полезно показать
  как «обратный фильм», закрепляя идею непрерывного семейства.
\end{remarkIT}
